\chapter{Introduction to quantum mechanics}

%%%%%%%%%%%%%%%%%%%%%%%%%%%%%
\section{Linear algebra}
Linear algebra is the study of vector spaces and of linear operations on those vector spaces.
The vector space of most interest for quantum mechanics is $\mathbb{C}^n$,
the set of all $n$-tuples of complex numbers $(z_1, \ldots, z_n)$.
Vector spaces have an \emph{addition} operation and a \emph{scalar multiplication} operation.
A vector space contains a \emph{zero vector} denoted by $0$
(not to be confused with $\ket{0}$),
so that $\ket{v} + 0 = \ket{v}$ for any vector $\ket{v}$.

A \emph{vector subspace} of a vector space $V$ is a subset $W \subseteq V$ that is itself a vector space
(W must be closed under addition and scalar multiplication, and contain the zero vector).

\subsection{Bases and linear independence}
A \emph{spanning set} for a vector space is a set of vectors $\ket{v_1}, \ldots, \ket{v_n}$ such that any vector $\ket{v} \in V$ can be expressed as a linear combination of the vectors in the spanning set:
\[
    \ket{v} = \sum_{i=1}^n a_i \ket{v_i}
\]

A set of non-zero vectors $\ket{v_1}, \ldots, \ket{v_n}$ is \emph{linearly dependent} if there exists a set of complex numbers $a_1, \ldots, a_n$, not all zero, such that
\[
    \sum_{i=1}^n a_i \ket{v_i} = 0
\]

A set of vectors is \emph{linearly independent} if it is not linearly dependent.
Any two sets of linearly independent vectors that span the same vector space $V$ contain the same number of vectors.
Such sets are called \emph{bases} for the vector space $V$, and always exists.
\graffito{From now on, we only consider finite-dimensional vector spaces.}
The number of elements in the basis is called \emph{dimension} of $V$.

\subsection{Linear operators and matrices}
A \emph{linear operator} $A$ between two vector spaces $V$ and $W$ is any function $A: V \to W$ which is linear in its inputs:
\[
    A\left(\sum_{i=1}^n a_i \ket{v_i}\right) = \sum_{i=1}^n a_i A\left(\ket{v_i}\right) \equiv \sum_{i=1}^n a_i A\ket{v_i}
\]
A linear operator $A$ defined on $V$ maps $V$ to itself, i.e. $A: V \to V$.

The \emph{identity operator} $I$ is the linear operator that maps every vector to itself: $I\ket{v} = \ket{v}$ for all $\ket{v} \in V$.
The \num{0} operator maps all vectors to the zero vector: $0\ket{v} = 0$ for all $\ket{v} \in V$.

If $A\colon V \to W$ and $B\colon W \to X$ are linear operators, then the \emph{composition} $BA\colon V \to X$ is also a linear operator defined by $(BA)\left(\ket{v}\right) = B\left(A\left(\ket{v}\right)\right)$.
We usually write $BA\ket{v}$ instead of $(BA)\left(\ket{v}\right)$.

Linear operators have a \emph{matrix representation} in a given basis.

A matrix is a linear operator since
\[
    A\left(\sum_{i=1}^n a_i \ket{v_i}\right) = \sum_{i=1}^n a_i A\ket{v_i}
\]
is true where $A$ is a matrix and $\ket{v_i}$ are column vectors.
Conversely, a linear operator can be given a matrix representation.
If $A\\colon V \to W$ is a linear operator,
and $\{\ket{v_1}, \ldots, \ket{v_n}\}$ and $\{\ket{w_1}, \ldots, \ket{w_m}\}$ are bases for $V$ and $W$, respectively,
for each $j=1, \ldots, m$ there exist complex numbers $A_{ij}$ for each $i=1, \ldots, n$ such that
\[
    A\ket{v_j} = \sum_{i=1}^m A_{ij}\ket{w_i}
\]
The matrix whose entries are $A_{ij}$ is called the \emph{matrix representation} of the linear operator $A$.

The matrix representation of $A$ is completely equivalent to the linear operator $A$.

\subsection{The Pauli matrices}
The \emph{Pauli matrices} are a set of three $2 \times 2$ complex matrices which are Hermitian and unitary.
\[
    \sigma_0 = I = \begin{pmatrix}
        1 & 0 \\
        0 & 1
    \end{pmatrix}
\]
\[
    \sigma_1 = \sigma_x = X = \begin{pmatrix}
        0 & 1 \\
        1 & 0
    \end{pmatrix}
\]
\[
    \sigma_2 = \sigma_y = Y = \begin{pmatrix}
        0 & -i \\
        i & 0
    \end{pmatrix}
\]
\[
    \sigma_3 = \sigma_z = Z = \begin{pmatrix}
        1 & 0 \\
        0 & -1
    \end{pmatrix}
\]

\subsection{Inner products}
An \emph{inner product} on a vector space $V$ is a function $(\cdot, \cdot): V \times V \to \mathbb{C}$ that satisfies the following properties:
\begin{itemize}
    \item $(\cdot, \cdot)$ is linear in its second argument
        \[
            \left(\ket{v}, \sum_i \lambda_i \ket{w_i}\right) = \sum_i \lambda_i (\ket{v}, \ket{w_i})
        \]
    \item $(\cdot, \cdot)$ is conjugate symmetric
        \[
            (\ket{v}, \ket{w}) = (\ket{w}, \ket{v})^*
        \]
    \item $(\cdot, \cdot)$ is positive definite
        \[
            (\ket{v}, \ket{v}) \geq 0 \quad \text{and} \quad (\ket{v}, \ket{v}) = 0 \iff \ket{v} = 0
        \]
\end{itemize}

A vector space equipped with an inner product is called an \emph{inner product space}.
For finite-dimensional complex vector spaces, the inner product space is an \emph{Hilbert space}.

In quantum mechanics, the inner product is usually denoted by $\braket{v|w}$ instead of $(\ket{v}, \ket{w})$. 
$\ket{v}$ is the \emph{dual vector} of $\bra{v}$.
The dual is a linear operator from the inner product space $V$ to $\mathbb{C}$ defined by $\bra{v}(\ket{w}) = \braket{v|w} = (\ket{v}, \ket{w})$.

$\mathbb{C}^n$ has an inner product defined by
\[
    \left((y_1, \ldots, y_n), (z_1, \ldots, z_n)\right) = \sum_{i=1}^n y_i^* z_i = (y^*_1, \ldots, y^*_n)\begin{pmatrix}
        z_1 \\
        \vdots \\
        z_n
\end{pmatrix}
\]

Vectors $\ket{v}$ and $\ket{w}$ are \emph{orthogonal} if $\braket{v|w} = 0$.
The \emph{norm} of a vector $\ket{v}$ is defined as $\norm{\ket{v}} = \sqrt{\braket{v|v}}$.
A \emph{unit vector} is a vector $\ket{v}$ such that $\norm{\ket{v}} = 1$.
$\ket{v}/\norm{\ket{v}}$ is the \emph{normalized} form of $\ket{v}$.
A set $\ket{i}$ of vectors with index $i$ is \emph{othonormal} if $\braket{i|j} = \delta_{ij}$.

Given a basis set $\{\ket{w_1}, \ldots, \ket{w_n}\}$ for a vector space $V$ with an inner product,
the \emph{Gram-Schmidt orthogonalization} procedure produces an orthonormal basis set $\{\ket{v_1}, \ldots, \ket{v_n}\}$ for $V$.
Define $\ket{v_1} = \ket{w_1}/\norm{\ket{w_1}}$.
For $1\leq k \leq d-1$, define $\ket{v_{k+1}}$ inductively by
\[
    \ket{v_{k+1}} = \frac{\ket{w_{k+1}} - \sum_{i=1}^k \braket{v_i|w_{k+1}}\ket{v_i}}{\norm{\ket{w_{k+1}} - \sum_{i=1}^k \braket{v_i|w_{k+1}}\ket{v_i}}}
\]
Any finite dimensional vector space of dimension $d$ has an orthonormal basis.\graffito{From now on, when we speak of a matrix representation for a linear operator,
we mean a matrix representation with respect to orthonormal input and output bases.
If the input and output spaces are the same, then we assume that the input and output bases are the same,
unless stated otherwise.}

The inner product of two vectors is equal to the vector inner product between two matrix representations of the vectors,
with respect to the same an orthonormal basis. 
The dual vector $\bra{v}$ has a matrix representation that is the conjugate transpose of the matrix representation of $\ket{v}$.

Linear operators can be represented using the inner product.
If $\ket{v}$ is a vector in an inner product space $V$,
and $\ket{w}$ is a vector in an inner product space $W$,
then the linear operator $\ket{w}\bra{v}$ maps $V$ to $W$ and is defined by
\[
    (\ket{w}\bra{v})\ket{v'} = \braket{v|v'}\ket{w}
\]
This representation is called \emph{outer product} representation.

The outer product representation gives us the \emph{completness relation} for orthonormal basis sets:
\[
    \sum_i \ket{i}\bra{i} = I
\]

One application of the completeness relation is to give a means for representing any operator in the outer product representation.
Suppose $A\colon V \to W$ is a linear operator,
$\ket{v_i}$ is an orthonormal basis for $V$, and $\ket{w_j}$ is an orthonormal basis for $W$.
Then we can write
\[
    A = I_W A I_V = \sum_{j} \ket{w_j}\bra{w_j} A \sum_i \ket{v_i}\bra{v_i} = \sum_{i,j} \braket{w_j|A|v_i}\ket{w_j}\bra{v_i} 
\]
which is the outer product representation of $A$.
A has matrix elements $A_{ji} = \braket{w_j|A|v_i}$ with respect to the orthonormal bases $\{\ket{v_i}\}$ and $\{\ket{w_j}\}$.

The \emph{Cauchy-Schwarz inequality} states that for any vectors $\ket{v}$ and $\ket{w}$ in an inner product space,
\[
    \abs{\braket{v|w}}^2 \leq \braket{v|v}\braket{w|w}
\]
It is an important geometric fact about Hilbert spaces.

\subsection{Eigenvectors and eigenvalues}
